\section*{Capítulo \ref{ch:b1:stress-strain-rheology} --- Tensão, deformação e reologia}

\begin{solution}{exo:b1:stress-strain-rheology:1}
(a) $\sigma_v = 2700\times9.81\times10\,000 = \qty{2.65e8}{Pa} = \qty{264.9}{MPa}$.
(b) $P = 1000\times9.81\times10\,000 = \qty{98.1}{MPa}$. (c) $\sigma_v' =
264.9 - 98.1 = \qty{166.8}{MPa}$, e $\lambda = 98.1/264.9 = 0.370$, que é
$1000/2700$.
\end{solution}

\begin{solution}{exo:b1:stress-strain-rheology:2}
(a) $e = (13-10)/10 = 0.30$ e $S = 1.30$. (b) $e = (7-10)/10 = -0.30$, um
encurtamento de \qty{30}{\percent}, e $S = 0.70$. (c) O ângulo se fechou de
$\psi = \ang{30}$, e $\gamma = \tan\ang{30} = 0.577$.
\end{solution}

\begin{solution}{exo:b1:stress-strain-rheology:3}
(a) Centro \qty{70}{MPa}, raio \qty{30}{MPa}. (b) A maior tensão de cisalhamento
é o raio, \qty{30}{MPa}, nos planos a \ang{45} de $\sigma_1$ e
$\sigma_3$, que carregam a \omterm{def:b1:stress-strain-rheology:stress}{tensão normal} do centro, \qty{70}{MPa}. (c)
$2\theta = \ang{60}$: $\sigma_n = 70 + 30\cos\ang{60} = \qty{85}{MPa}$ e
$\tau = 30\sin\ang{60} = \qty{26.0}{MPa}$.
\end{solution}

\begin{solution}{exo:b1:stress-strain-rheology:4}
(a) $\ang{45} - \ang{30}/2 = \ang{30}$ com $\sigma_1$, de um lado e do outro dele, logo os dois planos formam \ang{60} entre si. (b) Com $c' = 20 + r$, a
condição do \cref{thm:b1:stress-strain-rheology:failure-angle} lê-se
$r = C\cos\phi + (20 + r)\sin\phi$, logo
\[
  r = \frac{C\cos\phi + 20\sin\phi}{1 - \sin\phi} = \frac{8.66 + 10}{0.5}
    = \qty{37.3}{MPa},
\]
e $\sigma_1' = 20 + 2r = \qty{94.6}{MPa}$.
\end{solution}

\begin{solution}{exo:b1:stress-strain-rheology:5}
(a) $c = 40$ e $R = \sqrt{20^2 + 15^2} = \qty{25}{MPa}$: $\sigma_1 =
\qty{65}{MPa}$ e $\sigma_3 = \qty{15}{MPa}$; $\tan2\theta_p = 30/40 = 0.75$,
$2\theta_p = \ang{36.87}$, $\theta_p = \ang{18.43}$. (b) Nesse ângulo
$\cos^2\theta = 0.9$, $\sin^2\theta = 0.1$ e $\sin\theta\cos\theta = 0.3$:
$\sigma_n = 60\times0.9 + 2\times15\times0.3 + 20\times0.1 = 54 + 9 + 2 =
\qty{65}{MPa}$, e $\tau = (20-60)\times0.3 + 15\times(0.9-0.1) = -12 + 12 =
0$. (c) $R = \qty{25}{MPa}$, nos planos a \ang{45} das direções principais.
\end{solution}

\begin{solution}{exo:b1:stress-strain-rheology:6}
(a) $S_{1,2} = \sqrt{1 + 1} \pm 1$: $S_1 = 2.414$ e $S_2 = 0.414$;
$\tan2\theta' = 2/2 = 1$, $\theta' = \ang{22.5}$. (b) $S_1S_2 =
(\sqrt2 + 1)(\sqrt2 - 1) = 1$, que é $\det F$: nenhuma mudança de área. (c) A
reta ao longo de $(1, 1)$ vai a $F(1, 1) = (1 + 2, 1) = (3, 1)$, a
$\arctan(1/3) = \ang{18.4}$ da direção de cisalhamento, com o \omterm{def:b1:stress-strain-rheology:strain}{alongamento}
$\sqrt{10}/\sqrt2 = 2.24$. Ela girou para além do eixo longo da elipse, a \ang{22.5}, em direção ao plano de cisalhamento; seu \omterm{def:b1:stress-strain-rheology:strain}{alongamento} é um pouco menor que
$S_1$, como deve ser, pois nenhuma linha se alonga mais que $S_1$.
\end{solution}

\begin{solution}{exo:b1:stress-strain-rheology:7}
(a) $S_1/S_2 = 4$ e $S_1S_2 = 1$ dão $S_1 = 2$ e $S_2 = 0.5$; pela
\cref{prop:b1:stress-strain-rheology:simple-shear}, $\gamma = S_1 - S_2 =
1.5$. (b) $\tan2\theta' = 2/1.5 = 1.333$, $2\theta' = \ang{53.1}$: os eixos longos devem jazer a \ang{26.6} das paredes. (c) $1.5\times200 = \qty{300}{m}$.
\end{solution}

\begin{solution}{exo:b1:stress-strain-rheology:8}
$\sigma_v = 2700\times9.81\times3000 = \qty{79.46}{MPa}$. (a) $P = \qty{29.43}{MPa}$,
$\sigma_n' = 70 - 29.43 = \qty{40.57}{MPa}$, e Byerlee permite
$0.85\times40.57 = \qty{34.5}{MPa}$, mais que os \qty{30}{MPa} aplicados: a
falha se mantém. (b) $P = 0.6\times79.46 = \qty{47.68}{MPa}$, $\sigma_n' =
\qty{22.32}{MPa}$, e o atrito dá apenas \qty{19.0}{MPa}: a falha
desliza. (c) $30 = 0.85\,(70 - 79.46\,\lambda)$ dá $\lambda = (70 -
35.29)/79.46 = 0.437$.
\end{solution}

\begin{solution}{exo:b1:stress-strain-rheology:9}
(a) A tensão varia como $\dot\varepsilon^{1/4}$: $10^{1/4} = 1.78$. (b) A
razão das resistências é
\[
  \exp\Bigl[\frac{E}{nR}\Bigl(\frac{1}{673.15} - \frac{1}{773.15}\Bigr)\Bigr]
  = \exp\bigl(6705\times\num{1.921e-4}\bigr) = \exp(1.288) = 3.63 .
\]
(c) Uma variação de dez vezes na taxa muda a resistência à fluência por um fator 1.78, o que a temperatura faz em apenas $100\times\ln1.78/\ln3.63 \approx
\qty{45}{\degreeCelsius}$; com uma geoterma de cerca de \qty{20}{\degreeCelsius}
por quilômetro, isso desloca a transição em cerca de \qty{2}{km}. Uma mudança
na geoterma desloca as isotermas, e com elas a transição, muito mais.
\end{solution}

\begin{solution}{exo:b1:stress-strain-rheology:10}
(a) Com $C = 0$, uma falha desliza quando $r = \mu c'/\sqrt{1+\mu^2}$
(\cref{thm:b1:stress-strain-rheology:failure-angle}). Aqui $c' = S =
\rho g z(1-\lambda)$, logo $\sigma_1 - \sigma_3 = 2r = 2\mu S/\sqrt{1+\mu^2}$.
(b) $S$ cresce em $2700\times9.81\times1000\times(1 - 1000/2700) =
\qty{16.68}{MPa}$ por quilômetro, e $2\times0.85/1.3124 = 1.295$: a inclinação
é \qty{21.6}{MPa} por quilômetro, entre \num{13.1} (falhas normais) e
\qty{61.3}{MPa} por quilômetro (empurrões). (c) $2\mu/(k+\mu) = 0.786$,
$2\mu/k = 1.295$ e $2\mu/(k-\mu) = 3.68$, com $k = 1.3124$: a \omterm{def:b1:stress-strain-rheology:principal}{tensão diferencial} no deslizamento é 0.79, 1.30 e 3.68 vezes a \omterm{def:b1:stress-strain-rheology:lithostatic}{tensão efetiva} vertical.
\end{solution}

\begin{solution}{exo:b1:stress-strain-rheology:11}
(a) $t = 60\times\num{e6}\times\num{3.156e7} = \qty{1.89e15}{s}$ e
$2\sqrt{\kappa t} = 2\sqrt{\num{1.89e9}} = \qty{87.0}{km}$. (b)
$\operatorname{erf}x = 600/1350 = 0.444$ para $x = 0.417$: $z = 0.417\times87.0
= \qty{36.3}{km}$; $\operatorname{erf}x = 800/1350 = 0.593$ para $x = 0.586$:
$z = \qty{51.0}{km}$. (c) Os terremotos deveriam atingir \qtyrange{36}{51}{km}; a envoltória põe a base da camada frágil a \qty{38.2}{km}, dentro dessa faixa e perto de seu topo, a uma temperatura, \qty{641}{\degreeCelsius}, na faixa observada.
\end{solution}

\begin{solution}{exo:b1:stress-strain-rheology:12}
(a) $r = 45$ e $c' = \qty{65}{MPa}$; a distância do centro à reta
$\tau = 0.85\sigma_n'$ é $0.85\times65/1.3124 = \qty{42.1}{MPa}$, menor que o raio: o círculo cruza a reta, e os planos no arco acima dela deslizam. (b) Com $2\theta$ o ângulo a partir da extremidade $\sigma_1$, o deslizamento exige
$45\sin2\theta \ge 0.85\,(65 + 45\cos2\theta)$, isto é
$45\sin2\theta - 38.25\cos2\theta \ge 55.25$, ou
$59.06\sin(2\theta - \ang{40.36}) \ge 55.25$, onde $\ang{40.36} = \phi$.
Logo $\sin(2\theta - \phi) \ge 0.9355$, $2\theta - \phi$ entre \ang{69.31}
e \ang{110.69}, $2\theta$ entre \ang{109.67} e \ang{151.06}: as normais jazem a \ang{54.8} a \ang{75.5} de $\sigma_1$, e as fraturas que deslizam a \ang{14.5} a \ang{35.2} de $\sigma_1$. O meio da faixa,
\ang{24.8}, é $\ang{45} - \phi/2$, o plano do
\cref{thm:b1:stress-strain-rheology:failure-angle}. (c) Apenas esse plano desliza
quando o círculo apenas toca a reta: $r = \mu(20 + r)/k$, $r = 20\mu/(k -
\mu) = 17/0.4624 = \qty{36.8}{MPa}$, e $\sigma_1' = 20 + 2r =
\qty{93.5}{MPa}$.
\end{solution}

\begin{solution}{pb:b1:stress-strain-rheology:1}
\textbf{1.} $\sigma_v = 2700\times9.81\times5000 = \qty{132.44}{MPa}$.

\textbf{2.} $P = 1000\times9.81\times5000 = \qty{49.05}{MPa}$, e
$\lambda = 49.05/132.44 = 0.370$.

\textbf{3.} $\sigma_1 = \sigma_v = \qty{132.44}{MPa}$, vertical, e
$\sigma_3 = \qty{75.00}{MPa}$, horizontal leste--oeste: a \omterm{def:b1:layered-earth:layers}{crosta} é distendida leste--oeste sob o próprio peso, o estado que impulsiona falhas normais, como esta.

\textbf{4.} $\sigma_1 - \sigma_3 = \qty{57.44}{MPa}$; centro
$(132.44 + 75.00)/2 = \qty{103.72}{MPa}$; raio \qty{28.72}{MPa}.

\textbf{5.} $\sigma_1' = 132.44 - 49.05 = \qty{83.39}{MPa}$, $\sigma_3' =
75.00 - 49.05 = \qty{25.95}{MPa}$, centro $103.72 - 49.05 = \qty{54.67}{MPa}$.

\textbf{6.} $\sigma_{xx} = \qty{75.00}{MPa}$, $\sigma_{yy} = \qty{132.44}{MPa}$
e $\sigma_{xy} = 0$: os eixos são as direções principais, uma delas
vertical, de modo que os planos perpendiculares a eles não carregam tensão de cisalhamento.

\textbf{7.} Ao longo do \omterm{def:b1:geological-maps:attitude}{mergulho}, a falha segue o vetor unitário
$(\cos\ang{60}, -\sin\ang{60})$, para leste e para baixo; sua normal ascendente é
$\vec n = (\sin\ang{60}, \cos\ang{60}) = (0.866, 0.500)$, a \ang{30} de
$x$. Seu ângulo com $\sigma_1$, ao longo de $y$, tem cosseno
0.500: $\theta = \ang{60}$.

\textbf{8.} Com \ang{30} de $x$: $\sigma_n = 75.00\times0.75 + 132.44\times0.25
= 56.25 + 33.11 = \qty{89.36}{MPa}$, e $\tau = (132.44 - 75.00)\times
\sin\ang{30}\cos\ang{30} = 57.44\times0.4330 = \qty{24.87}{MPa}$.

\textbf{9.} Em $2\theta = \ang{120}$ a partir da extremidade $\sigma_1$:
$\sigma_n = 103.72 + 28.72\cos\ang{120} = 103.72 - 14.36 = \qty{89.36}{MPa}$
e $\tau = 28.72\sin\ang{120} = \qty{24.87}{MPa}$.

\textbf{10.} $\sigma_n' = 89.36 - 49.05 = \qty{40.31}{MPa}$.

\textbf{11.} $24.87/40.31 = 0.617$, abaixo de 0.85: a falha se mantém.

\textbf{12.} $0.85\times40.31 = \qty{34.26}{MPa}$; a tensão de cisalhamento fica aquém dela em \qty{9.39}{MPa}.

\textbf{13.} O círculo tem centro em \qty{54.67}{MPa} e raio
\qty{28.72}{MPa}, de \num{25.95} a \qty{83.39}{MPa}; a reta de Byerlee é
$\tau = 0.85\,\sigma_n'$ (todas as tensões normais aqui estão abaixo de
\qty{200}{MPa}); a falha é o ponto $(40.31, 24.87)$, a \ang{120} da
extremidade $\sigma_1$.

\textbf{14.} $0.85\times54.67/\sqrt{1 + 0.85^2} = 46.47/1.3124 =
\qty{35.41}{MPa}$, mais que o raio, \qty{28.72}{MPa}: o círculo inteiro
fica abaixo da reta, e nenhum plano através da rocha pode deslizar.

\textbf{15.} $\phi = \arctan0.85 = \ang{40.36}$: o plano a
$\ang{45} - \phi/2 = \ang{24.8}$ do $\sigma_1$ vertical, mergulhando
\ang{65.2}.

\textbf{16.} $S = \sigma_v' = \qty{83.39}{MPa}$ e
$\sigma_1 - \sigma_3 = 2\times0.85\times83.39/(1.3124 + 0.85) =
\qty{65.55}{MPa}$; a inclinação dá $13.11\times5 = \qty{65.55}{MPa}$.

\textbf{17.} $\sigma_3 = 132.44 - 65.55 = \qty{66.88}{MPa}$.

\textbf{18.} A falha desliza quando
$r\sin\ang{120} = 0.85\,(c' + r\cos\ang{120})$ com
$c' = (83.39 + \sigma_3')/2$ e $r = (83.39 - \sigma_3')/2$; resolvendo obtemos
$\sigma_3' = \qty{17.18}{MPa}$, $\sigma_3 = \qty{66.23}{MPa}$. Ela é menor
do que para o melhor plano: a falha não jaz no ponto em que o círculo toca a reta pela primeira vez, de modo que o círculo deve crescer mais antes de atingir o próprio ponto da falha.

\textbf{19.} Toda \omterm{def:b1:stress-strain-rheology:stress}{tensão normal} efetiva, $\sigma_1'$ e $\sigma_3'$ incluídas, cai de $\Delta P$, e sua diferença não muda; as tensões de cisalhamento não dependem da pressão nos poros de modo algum.

\textbf{20.} O melhor plano desliza quando a distância é igual ao raio:
$0.85\,(54.67 - \Delta P)/1.3124 = 28.72$, logo
$\Delta P = 54.67 - 28.72\times1.3124/0.85 = 54.67 - 44.34 = \qty{10.33}{MPa}$.

\textbf{21.} A coluna de água no poço pesa o mesmo de antes, de modo que a
bomba deve acrescentar toda a \qty{10.33}{MPa} na boca do poço, cerca de 103 vezes
a pressão atmosférica: o peso de uma coluna de água
$\num{10.33e6}/(1000\times9.81) = \qty{1053}{m}$ de altura.

\textbf{22.} Com $\mu = 0.6$ a falha poderia suportar apenas
$0.6\times40.31 = \qty{24.19}{MPa}$, menos que os \qty{24.87}{MPa} que ela carrega: ela já estaria deslizando. Como está quieta, seu atrito é pelo menos $0.617$, ou ela tem alguma \omterm{def:b1:stress-strain-rheology:coulomb}{coesão}.

\textbf{23.} O plano mais bem orientado desliza quando a pressão nos poros atinge
$49.05 + 10.33 = \qty{59.38}{MPa}$, isto é, $\lambda = 59.38/132.44 = 0.448$,
bem abaixo do valor litostático, 1.

\textbf{24.} $\Delta P = \sigma_n' - \tau/0.85 = 40.309 - 24.870/0.85 = 40.309 -
29.259 = \qty{11.049}{MPa}$, mais que os \qty{10.33}{MPa} que moveriam um plano da melhor orientação. Com uma casa decimal, a elevação da pressão nos poros que leva a falha a deslizar é
\textbf{\qty[text-series-to-math=true]{11.0}{MPa}}.
\end{solution}
